% =============================================================
%  dissertation_preamble.tex
%  Shared package imports and custom commands.
%  \input this file at the top of dissertation.tex ONLY.
% =============================================================

% --- Page geometry -----------------------------------------------------------
% NOTE: caltech_thesis.cls already loads `geometry` itself (1.5in side
% margins, 1in top/bottom, per Caltech's binding/margin requirements), so
% this preamble must NOT load it again with different options -- LaTeX
% would apply whichever load wins and silently ignore the other, which is
% exactly the kind of margin mismatch that gets a thesis bounced at
% submission. If dissertation.tex is ever compiled under a different
% (non-Caltech) document class that does NOT already set geometry, restore
% a line like `\usepackage[margin=1in, letterpaper]{geometry}` here.

% --- Typography --------------------------------------------------------------
\usepackage{lmodern}          % scalable Latin Modern fonts
\usepackage[T1]{fontenc}
\usepackage[utf8]{inputenc}
\usepackage{microtype}        % improved justification
\usepackage{textcomp}         % \textperthousand (per mille, U+2030) below --
                               % most other common unicode punctuation (curly
                               % quotes, en/em dash, degree, plus-minus,
                               % ellipsis) is already handled by inputenc's
                               % utf8 definitions and needs nothing extra.

% A few characters straight from pasted/pandoc-converted source text (e.g.
% Chapter 5's field-report content, which uses literal Greek delta for
% isotope notation like the literal characters "delta-18-O" rather than $\delta^{18}$O) aren't
% in inputenc's default utf8 active-character table and would otherwise be a
% fatal "Unicode character ... not set up for use with LaTeX" error. Map them
% here once, globally, rather than hand-converting every instance in the
% source text (safer -- avoids silently mistyping a value while
% converting).
\usepackage{newunicodechar}
\newunicodechar{δ}{\ensuremath{\delta}}
\newunicodechar{Δ}{\ensuremath{\Delta}}
\newunicodechar{‰}{\textperthousand}

% --- Math --------------------------------------------------------------------
\usepackage{amsmath, amssymb, amsthm}
\usepackage{bm}               % \bm bold math -- used by Chapter 4 (nGibbs methods)
\usepackage{siunitx}          % consistent SI units, e.g. \SI{4.2}{\mega\year}

% --- Figures & tables --------------------------------------------------------
\usepackage{graphicx}
\usepackage{calc}             % arithmetic in lengths (\real{}, +, -, *) -- needed
                               % by Pandoc-generated longtable column specs like
                               % p{(\linewidth - 12\tabcolsep) * \real{0.1951}},
                               % used by Chapter 3's Table 1 (Chapter3/Writing/
                               % HeInSulfides_Main_Body.tex)
\usepackage{float}
\usepackage{placeins}         % \FloatBarrier -- forces pending floats to
                               % resolve before a fixed point (e.g. right
                               % before a non-floating longtable/supertabular
                               % caption) without pinning them to their
                               % source position via [H], which can force a
                               % float onto a page it doesn't fit on and run
                               % it off the bottom margin. See chapter1_SI_
                               % body.tex (fig:predicted-vs-measured-permeation)
                               % for the case this was added for.
\usepackage{booktabs}         % professional table rules
\usepackage{longtable}        % tables that span pages
\usepackage{supertabular}     % multi-page table in Chapter 1 (apatite step-heat
                               % data), shared verbatim with TubeMethodPaper.tex,
                               % which loads supertabular natively via copernicus.cls
\usepackage{multirow}
\usepackage{caption}
\usepackage{subcaption}
\usepackage{array}           % extended column-spec features (>{...},
                              % \newcolumntype); harmless to keep loaded
                              % even though no table currently needs it
                              % now that Chapter 1's widest tables moved
                              % to \chplandscape (landscape) instead of
                              % font-shrinking -- see chapter1_body.tex.
\usepackage{pdflscape}        % \landscape environment, used via the
                              % \chplandscape wrapper (defined per-wrapper
                              % in Chapter1/chapter1.tex and Chapter1/
                              % chapter1_SI.tex) to rotate Chapter 1's
                              % widest tables for extra width in the
                              % thesis only -- the standalone
                              % TubeMethodPaper.tex/Paper/SI.tex define
                              % \chplandscape as a no-op instead, keeping
                              % those journal manuscripts portrait.

% Auto-center every table and shrink its text, without editing individual
% tables. Caltech's required 1.5in side margins leave only a 5.5in text
% width -- narrower than most of these tables were originally drafted for,
% which is why they were rolling off the right edge while left-justified.
%
% NOTE: an earlier version of this tried to wrap every \tabular in an
% adjustbox (\AtBeginEnvironment{tabular}{\begin{adjustbox}{max
% width=\textwidth}}) to hard-guarantee a fit regardless of native column
% widths. That broke the build (fatal "Missing }"/"invalid character"
% errors) -- \AtBeginEnvironment cannot safely wrap an environment that
% takes a mandatory brace argument (tabular's {column-spec}) in another
% such environment, because the injected \begin{adjustbox}{...} steals the
% argument-parsing slot tabular expects for its own column spec. Plain
% declarations (\centering, \footnotesize, \scriptsize -- no arguments of
% their own) are safe here; wrapping one environment inside another via
% this mechanism is not.
\AtBeginEnvironment{table}{\centering\footnotesize}
\AtBeginEnvironment{longtable}{\centering\scriptsize}
\AtBeginEnvironment{supertabular}{\centering\scriptsize}
\setlength{\tabcolsep}{3pt}   % default is 6pt -- tighter padding buys back
                              % real width on many-column numeric tables

% A few of the widest tables use fixed-width p{...} columns (in cm) rather
% than content-driven columns, so shrinking the font alone doesn't reclaim
% any width -- those still need a per-table \resizebox wrap (done directly
% in the affected .tex files, not here) for a hard guarantee of fit.
\graphicspath{
  {figures/}
  {Chapter1/Figures/source/}
  {Chapter1/Figures/generated/}
  {Chapter2/Figures/source/}
  {Chapter2/Figures/generated/}
  {Chapter3/Figures/source/}
  {Chapter3/Figures/generated/}
  {Chapter3/Figures/generated/step_heat_regressions/}
  {Chapter4/Figures/}
  {Chapter4/Figures/generated/}
  {Chapter5/Figures/}
  {Chapter5/Figures/source/}
  {Chapter5/Figures/generated/}
  {ChapterA/Figures/}
}

% --- Colors & hyperlinks -----------------------------------------------------
\usepackage[dvipsnames]{xcolor}
\usepackage{colortbl}         % \cellcolor shading in the Chapter 4 quality tables
\usepackage[
  colorlinks = true,
  linkcolor  = NavyBlue,
  citecolor  = OliveGreen,
  urlcolor   = Maroon,
  pdftitle   = {From Surface to Core},
  pdfauthor  = {Ben Thyer},
]{hyperref}
\usepackage{doi}

% --- Bibliography ------------------------------------------------------------
% Name-truncation rule (discipline convention, not a Caltech requirement --
% Caltech's thesis guidelines explicitly leave reference/citation style to
% "the style current in the thesis writer's discipline," only requiring it be
% used consistently throughout): 2-author works print both names; 3+-author
% works truncate to the first author + "et al." (maxnames=2, minnames=1).
% uniquelist=true is the standard biblatex safety net for the case where two
% DIFFERENT works would otherwise truncate to an identical-looking citation
% (e.g. "Smith et al., 2020" for two unrelated Smith-led 2020 papers) -- it
% expands the printed name list just enough to tell them apart (e.g. "Smith,
% Jones, et al." vs. "Smith, Brown, et al.") rather than appending a/b/c
% letters. True letter disambiguation (2020a / 2020b) already happens
% automatically -- with no extra config -- whenever two entries share the
% exact same full author list AND year (e.g. a genuine same-author same-year
% duplicate); it does not, by default, cover the "different authors, same
% truncated display" case, which is what uniquelist handles instead. As of
% this writing, a scan of references.bib found zero actual collisions under
% the new maxnames=2 rule, so this is a forward-looking safety net, not a
% fix for a live problem.
%
% maxbibnames is set high (effectively "no truncation") so the References
% list at the end always spells out every author, regardless of how the
% in-text citation truncates -- maxnames/minnames above only govern in-text
% \citep/\citet, not the bibliography.
\usepackage[
  backend    = biber,
  style      = authoryear,   % change to e.g. agu or geology if preferred
  sorting    = nyt,
  maxnames   = 2,
  minnames   = 1,
  maxbibnames = 99,
  uniquelist = true,
  giveninits = true,
  uniquename = init,
  natbib     = true,         % provides \citep/\citet/\citealp (natbib syntax),
                              % which Chapter 1's shared body (reused from
                              % TubeMethodPaper.tex) uses throughout. Without
                              % this, biblatex only offers its own
                              % \parencite/\textcite commands.
]{biblatex}
\addbibresource{references.bib}   % master .bib file — keep alongside dissertation.tex

% --- Header / footer ---------------------------------------------------------
\usepackage{fancyhdr}
\pagestyle{fancy}
\fancyhf{}
\fancyhead[LE,RO]{\thepage}
\fancyhead[RE]{\textit{\leftmark}}
\fancyhead[LO]{\textit{\rightmark}}
\renewcommand{\headrulewidth}{0.4pt}

% caltech_thesis.cls sets headsep to a very tight ~0.25in - 1em (roughly
% 6pt), tuned for its own default page style, which puts only a bare page
% number in the header. Since fancyhdr above instead puts the italicized
% chapter/section title (\leftmark/\rightmark) in the header, long titles
% that wrap to a second line collide with the body text below. Calling
% \geometry{} again (rather than re-\usepackage-ing it) only updates the
% keys given here and safely recomputes \textheight to match -- it does
% NOT reload/conflict with the page-margin setup the NOTE above warns
% about. headheight=30pt comfortably fits two lines of a 12pt header;
% headsep=18pt keeps a clean gap below it. If a chapter title ever needs
% a third header line, bump headheight further (e.g. to 42pt).
\geometry{headheight=30pt, headsep=18pt}

% --- Chapter / section formatting --------------------------------------------
% NOTE: caltech_thesis.cls (memoir-based) already defines its own official
% Caltech chapter styling -- centered, italic small-caps "Chapter I" numbering
% in Roman numerals via \chapnamefont/\chaptitlefont/\printchapterheading -- so
% titlesec's \titleformat{\chapter} override below has been disabled. Loading
% titlesec's chapter redefinition on top of memoir's own chapter-heading system
% is a known source of conflicts/incompatibility warnings. Re-enable this only
% if compiling Chapter files under a plain `book`/`report` class instead.
% \usepackage{titlesec}
% \titleformat{\chapter}[display]
%   {\normalfont\Large\bfseries}{\chaptertitlename\ \thechapter}{20pt}{\Huge}
% \titlespacing*{\chapter}{0pt}{40pt}{20pt}

% --- Per-chapter Supplementary Information (SI) -------------------------------
% Each chapter's SI (Chapter<N>/chapter<N>_SI.tex) is \input{} immediately
% after that chapter's main body, inside the SAME \chapter{} -- it is not a
% chapter of its own. Two problems that creates, and how these commands fix
% them:
%
% 1. Visually, the SI used to just continue as the next plain \section{},
%    indistinguishable from the chapter's own sections. \SItitlepage{...}
%    gives it a proper divider/title page first (unnumbered, no header, own
%    page), the way a book's appendix or backmatter would get one.
%
% 2. Numerically, its \section headings used to continue the chapter's own
%    numbering (...7, 8, 9...), so a reader landing mid-SI from a link or a
%    flipped-to page couldn't tell from the running header alone whether they
%    were still in the chapter proper or in its supplement. The `ChapterSI`
%    environment resets the section counter and switches \thesection to a
%    "<chapter>.<n>s" format for everything between \begin{ChapterSI} and
%    \end{ChapterSI} (e.g. the SI's first section becomes "1.1s" instead of
%    continuing on from wherever the main text left off). \thesubsection
%    inherits this automatically, since its default definition already
%    builds on \thesection. \S\ref{}/running headers/the ToC all pick this
%    up for free, because \setcounter and \renewcommand are both local by
%    default, and LaTeX environments are always wrapped in their own group,
%    so everything reverts automatically at \end{ChapterSI} with no matching
%    "restore" step to maintain. (The *next* chapter's own section numbering
%    isn't at risk even without that revert -- \chapter{} always resets
%    \c@section to 0 globally on its own, regardless of whatever it was left
%    at.)
%
%    NOTE on naming, two dead ends hit along the way, for the record:
%     (a) This is a `\newenvironment`, not a pair of \newcommand{\beginSI}/
%         \newcommand{\endSI} macros -- LaTeX reserves all command names of
%         the form \end<word> for environment closers, so
%         \newcommand{\endSI}{...} is flatly illegal ("Or name \end...
%         illegal", LaTeX2e kernel restriction, independent of memoir/
%         caltech_thesis.cls). Defining it as a proper environment sidesteps
%         that and is the more idiomatic way to express "everything between
%         these two markers" in LaTeX anyway.
%     (b) The environment is named `ChapterSI`, not `SI` -- \newenvironment{SI}
%         would auto-generate \SI as its begin-code control sequence, which
%         collides with siunitx's own \SI{value}{unit} command (loaded
%         above, e.g. \SI{4.2}{\mega\year}). `ChapterSI` generates
%         \ChapterSI/\endChapterSI instead, which don't collide with
%         anything.
%    Repeat this same \SItitlepage{...}/\begin{ChapterSI}...\end{ChapterSI}
%    pattern for Chapter 2 and beyond's SI wrappers.
\newcommand{\SItitlepage}[1]{%
  \clearpage
  \thispagestyle{empty}
  \vspace*{\fill}
  \begin{center}
    {\Huge\bfseries Supplementary Information}\\[1.5em]
    {\Large for Chapter \thechapter: \itshape #1}
  \end{center}
  \vspace*{\fill}
  \clearpage
}
\newenvironment{ChapterSI}{%
  \setcounter{section}{0}
  \renewcommand{\thesection}{\thechapter.\arabic{section}s}
}{}

% --- Appendix ----------------------------------------------------------------
% NOT loaded explicitly: memoir (which caltech_thesis.cls is built on) already
% emulates the `appendix` package internally with its own default options, so
% an explicit `\usepackage[toc, page]{appendix}` here causes a fatal "Option
% clash for package appendix" error. The `\begin{appendices}...\end{appendices}`
% environment used in dissertation.tex still works via memoir's built-in
% emulation without loading the package separately.

% --- Miscellaneous -----------------------------------------------------------
\usepackage{enumitem}
% pdflscape already loaded above (Figures & tables section) for \chplandscape
\usepackage{pdfpages}         % include external PDFs (field maps, etc.)
\usepackage{soul}             % \hl{} highlighting -- used once in Chapter 2's
                               % shared body (chapter2_body.tex)
\usepackage{lineno}           % line numbers — comment out for final version
% \linenumbers

% --- Code listings with cached output ------------------------------------------
% A script is shown verbatim with \lstinputlisting[style=thesispython], followed
% by what it printed (\cachedoutput) and the figure(s) it drew. Results are
% never computed by LaTeX: Chapter4/Figures/run_cached_listing.py runs the script
% once and writes generated/<stem>/{prints.txt,fig1.pdf,...}, rerunning only when
% the script's source changes (`make listings`). Paths are relative to the
% dissertation root.
\usepackage{listings}
\lstdefinestyle{thesispython}{
  language=Python,
  basicstyle=\ttfamily\footnotesize,
  keywordstyle=\color{NavyBlue}\bfseries,
  commentstyle=\color{OliveGreen}\itshape,
  stringstyle=\color{BrickRed},
  numbers=left, numberstyle=\tiny\color{gray}, numbersep=6pt,
  frame=single, rulecolor=\color{gray!60},
  breaklines=true, columns=fullflexible, keepspaces=true,
  showstringspaces=false, upquote=true,
}
\lstdefinestyle{thesisoutput}{
  basicstyle=\ttfamily\footnotesize,
  frame=leftline, framerule=1.5pt, rulecolor=\color{gray!60},
  backgroundcolor=\color{gray!8}, xleftmargin=1em,
  breaklines=true, columns=fullflexible, keepspaces=true,
}
% \cachedoutput{<path to prints.txt>} -- placeholder box instead of a fatal
% error if the cache hasn't been generated yet (e.g. fresh clone).
\newcommand{\cachedoutput}[1]{%
  \IfFileExists{#1}%
    {\lstinputlisting[style=thesisoutput]{#1}}%
    {\fbox{\footnotesize Missing cached output \texttt{\detokenize{#1}} -- run \texttt{make listings}}}}

% --- Custom commands ---------------------------------------------------------
\newcommand{\degC}{\ensuremath{^\circ\mathrm{C}}}
\newcommand{\He}{\textsuperscript{4}He}
\newcommand{\UThHe}{U--Th/He}
\newcommand{\fO}{\ensuremath{f_{\mathrm{O_2}}}}
\newcommand{\fS}{\ensuremath{f_{\mathrm{S_2}}}}
\newcommand{\fH}{\ensuremath{f_{\mathrm{H_2O}}}}
\newcommand{\Ea}{\ensuremath{E_a}}
\newcommand{\todo}[1]{\textcolor{red}{\textbf{[TODO: #1]}}}
\providecommand{\code}[1]{\texttt{\detokenize{#1}}}
